\documentclass[11pt]{article}
\usepackage[a4paper,margin=23mm]{geometry}
\usepackage{amsmath,amssymb,amsthm,mathtools,bm,booktabs,microtype,xurl,hyperref}
\hypersetup{colorlinks=true,linkcolor=blue,urlcolor=blue,citecolor=blue,
pdftitle={P54 same-prescription gauge momentum and exact action certificate}}
\setlength{\emergencystretch}{2em}
\newcommand{\tr}{\operatorname{tr}}
\newcommand{\Tr}{\operatorname{Tr}}
\newcommand{\diag}{\operatorname{diag}}
\newcommand{\MS}{\overline{\mathrm{MS}}}
\newcommand{\FP}{\operatorname{FP}}
\title{Route-F P54: Same-Prescription Gauge Momentum,\\
Mixed Loops and Exact Transverse Action Identities}
\author{P54PQ-v2 four-dimensional mainline}
\date{17 September 2026}
\begin{document}
\maketitle
\begin{abstract}
We complete the one-loop Euclidean radial bosonic background kernel for
an explicit background-field gauge fixing of the unchanged P54 action.
Vector, scalar--vector and direct ghost contributions are reduced in
$d=4-2\epsilon$ before subtraction; finite rational terms are retained.
Zero-momentum vector CW and tadpole limits reproduce the existing common
prescription, while nonzero momentum retains all scalar internal states.
An independent Gaussian-integer calculation certifies the rank-50 and
rank-36 transverse action identities, including all mixed covectors.
No quartic scan, parameter replacement or physical flavor fit is made.
The result is not a Nielsen-certified physical pole or an all-orders
stability theorem. The large local hard-sector insertion remains.
\end{abstract}

\section{Frozen input and meaning of completion}
The potential is precisely \texttt{potential\_factory} in
\texttt{verify\_p54\_p1\_hessian\_spectrum.py}, with 328 canonical real
scalar coordinates and the original real parameter card. We use
\begin{align}
 X&=Rr,\quad r=(w,\sigma,v_s)=(1,.1265,.25)\omega,\\
 G&=R^TR=\diag(12/5,2,1),\quad
 g=\mu/\omega=.5626028899853225.
\end{align}
All displayed numerical matrices use $\omega=1$. Define
$H=V_0''(X)$, $T_r=\partial_r H$, $U_{rs}=\partial_r\partial_s H$,
with the potential parameters fixed during these derivatives.
The 38 soft scalar modes comprise 33 gauge Goldstones, one PQ mode and
four real tuned doublet components; there are 290 positive scalar modes.
Of the 45 gauge vectors, 33 are massive and 12 are unbroken.

\emph{Completion} below means the three-by-three radial restriction of the
one-loop \emph{bosonic background functional}, at nonzero Euclidean
$s=p_E^2$, for the explicit fixing in the next section. It does not mean
the full neutral physical mixing block, all nonradial amplitudes, or
all Wilson coefficients. The three unfitted radial Majorana masses
remain an explicit fermion input. The common invariant mass counterterms
are unchanged; no finite kinetic counterterm is tuned to cancel a result.

\section{Quadratic operator from the kinetic action}
Let $T^{aT}=-T^a$ be the actual real scalar representation matrices and
choose
\begin{equation}
 D_\mu X=\partial_\mu X+g a_\mu^aT^aX,\quad X=\bar X+\eta,\quad
 v_a=T^a\bar X,\quad
 F^a=\partial_\mu a_\mu^a+\xi g v_a^T\eta,\quad
 {\cal L}_{\rm gf}=\frac{F^aF^a}{2\xi}.
\end{equation}
The gauge background is zero. The mixed kinetic term before gauge fixing is
\[
 g a_\mu^a\{v_a^T\partial_\mu\eta
                -(\partial_\mu v_a)^T\eta\}.
\]
Integrating $g(\partial_\mu a_\mu^a)v_a^T\eta$ by parts cancels
the first term and adds a second copy of the latter. Consequently
\begin{align}
 {\cal Q}_B&=\begin{pmatrix}Q_{SS}&W^T\\W&Q_{VV}\end{pmatrix},
 &W^a_{\mu I}&=-2g(T^a\partial_\mu\bar X)_I,\\
 Q_{SS}&=-\partial^2+H(\bar X)+\xi g^2vv^T,&
 Q_c&=-\partial^2+\xi M^2,\\
 (Q_{VV})^{ab}_{\mu\nu}
 &=\left[-\partial^2\delta_{\mu\nu}
 +(1-\xi^{-1})\partial_\mu\partial_\nu\right]\delta^{ab}
 +M^2_{ab}\delta_{\mu\nu},&
 M^2&=g^2v^Tv.
\end{align}
Here $v$ is the $328\times45$ orbit matrix. The constant-background
mixed operator vanishes, but its first radial derivative at momentum
$p$ is $W_r=-2igp_\mu T^aR_r$. This is the term that cannot be inferred
from a constant-background vector mass spectrum.

The one-loop bosonic functional and its second derivative are
\begin{align}
 \Gamma_1^B&=\tfrac12\Tr\log{\cal Q}_B-\Tr\log Q_c+\Gamma_{\rm CT},\\
 \Pi^B_{rs}
 &=\tfrac12\Tr[{\cal Q}_B^{-1}{\cal Q}_{B,rs}
       -{\cal Q}_B^{-1}{\cal Q}_{B,r}{\cal Q}_B^{-1}{\cal Q}_{B,s}],
\end{align}
with the ghost expression having coefficient $-1$ instead of $1/2$.
This fixes the combinatorial factors and signs without assigning a
separate phenomenological projection.

\subsection{Actual vertices and finite gauge parameter}
Write $d_{ar}=T^aR_r$. Then
\begin{align}
 B_{r,ab}&=g^2(d_{ar}^Tv_b+v_a^Td_{br}),\\
 C_{rs,ab}&=g^2(d_{ar}^Td_{bs}+d_{as}^Td_{br}),\\
 H_\xi&=H+\xi g^2vv^T,\\
 T^\xi_r&=T_r+\xi g^2(d_rv^T+vd_r^T),\\
 U^\xi_{rs}&=U_{rs}+\xi g^2(d_rd_s^T+d_sd_r^T).
\end{align}
Rotate $B,C$ and $d$ by the same vector mass eigenbasis.
The unbroken 12-dimensional vector subspace is annihilated by these
radial vertices, as directly checked from the representation.
Hence the following masters only need strictly positive vector masses.
Scalar masses $b$ may vanish. At the stationary point $Hv=0$; finite
$\xi$ gives the 33 gauge Goldstones masses $\xi m_V^2$, leaving five
scalar zero modes. All finite-$\xi$ terms above must be retained together.

\section{Dimensional reduction of the momentum masters}
To avoid sign ambiguity, strip the overall $1/(16\pi^2)$ from every
integral and write
\begin{align}
 {\cal A}(a)&=a[\log(a/\mu^2)-1],\qquad {\cal A}(0)=0,\\
 {\cal B}(a,b;s)&=-\int_0^1 dx\,
       \log\frac{(1-x)a+xb+x(1-x)s}{\mu^2}.
\end{align}
The Euclidean tadpole has UV residue $-a$ and finite part $\mathcal A$;
the bubble has residue $+1$ and finite part $\mathcal B$.
The previous scalar code uses $L=-{\cal B}$. Useful limits are
\begin{align}
 {\cal B}(a,a;0)&=-\log(a/\mu^2),\\
 {\cal B}(a,0;0)&=1-\log(a/\mu^2),\\
 {\cal B}(0,0;s)&=2-\log(s/\mu^2),\quad s>0.
\end{align}
The doubly massless bubble at $s=0$ is not a finite local quantity.

\subsection{Tensor numerators}
Let $q=k+p$, $p^2=s$, and define finite parts
\begin{align}
 {\cal F}(a,b;s)
 &:=4\,\FP\!\int_k\frac{(k\cdot q)^2}{(k^2+a)(q^2+b)}
 \nonumber\\
 &=-(3a+b+s){\cal A}(a)-(a+3b+s){\cal A}(b)
       +(a+b+s)^2{\cal B}(a,b;s),\\
 {\cal J}(a,b;s)
 &:=\FP\!\int_k\frac{(p\cdot k)^2}{(k^2+a)(q^2+b)}
 \nonumber\\
 &=\frac14\{(a-b-s){\cal A}(a)+(b-a+3s){\cal A}(b)
       +(a-b-s)^2{\cal B}(a,b;s)\}.
\end{align}
The symbol $\int_k$ in these master definitions includes the stripped
$16\pi^2$ normalization. To derive them, substitute
$2k\cdot q=k^2+q^2-s$ and
$2p\cdot k=q^2-k^2-s$ before integrating. For example
\[
 \int_k\frac{k^2}{q^2+b}=(-b+s)\int_k\frac1{q^2+b},
\]
because after $k=q-p$ the odd term and the scaleless constant vanish in
dimensional regularization. Expanding the squared numerators gives the
displayed identities. Their UV residues are respectively
\begin{equation}
 F_{\rm pole}=4a^2+4b^2+4ab+3s(a+b)+s^2,\qquad
 J_{\rm pole}=\tfrac14s(s-a-b).
\end{equation}

\subsection{Vector and mixed masters}
For $a>0$, the vector propagator is
\begin{equation}
 D_{a,\mu\nu}(k)
 =\frac{P^T_{\mu\nu}(k)}{k^2+a}
       +\frac{\xi P^L_{\mu\nu}(k)}{k^2+\xi a}
 =\frac{\delta_{\mu\nu}}{k^2+a}
 +\frac{k_\mu k_\nu}{a}
   \left(\frac1{k^2+a}-\frac1{k^2+\xi a}\right).
\end{equation}
Set $c=\xi a$, $d=\xi b$. Its trace and two-propagator master are
\begin{align}
 {\cal A}_\xi^V(a)
 &=3{\cal A}(a)+2a+\xi{\cal A}(\xi a),\\
 {\cal T}_\xi(a,b;s)
 &=2{\cal B}(a,b;s)-2
       +\xi[{\cal B}(a,d;s)+{\cal B}(c,b;s)]\nonumber\\
 &\quad+\frac{{\cal F}(a,b;s)-{\cal F}(a,d;s)
                    -{\cal F}(c,b;s)+{\cal F}(c,d;s)}{4ab}.
\end{align}
In the trace of $D_aD_b$, the first three terms before subtraction
combine to $(d_{\rm spacetime}-2)B_{ab}+\xi B_{ad}+\xi B_{cb}$.
Therefore setting $d_{\rm spacetime}=4$ prematurely loses the finite
$-2$. Similarly $(3-2\epsilon)A_{\rm bare}$ yields $3{\cal A}+2a$.
The vector bubble UV pole is $3+\xi^2$; it has no momentum-dependent
UV pole. The seagull pole is $-(3+\xi^2)a$.

The scalar--vector master from two $W$ vertices is
\begin{equation}
 {\cal I}_\xi(a,b;s)
 :=\FP\!\int_k\frac{p_\mu D_{a,\mu\nu}(k)p_\nu}{q^2+b}
 =s{\cal B}(a,b;s)
       +\frac{{\cal J}(a,b;s)-{\cal J}(\xi a,b;s)}{a}.
 \label{eq:mixedmaster}
\end{equation}
Its UV residue is $s(3+\xi)/4$. Its value at zero momentum is zero,
but its derivative is not:
\begin{equation}
 {\cal I}'_\xi(a,b;0)
 =\frac34{\cal B}(a,b;0)
       +\frac{\xi}{4}{\cal B}(\xi a,b;0)+\frac{\xi-1}{8}.
 \label{eq:mixedslope}
\end{equation}
At zero momentum, angular averaging gives
$(1-1/d)B_{ab}+(\xi/d)B_{\xi a,b}$ before subtraction.
Expanding $1/d=1/4+\epsilon/8+\cdots$ proves the finite last term.
In particular, the Landau $-1/8$ is not adjustable.

Checks independent of the model include
\begin{align}
 {\cal T}_1(a,b;s)&=4{\cal B}(a,b;s)-2,&
 {\cal I}_1(a,b;s)&=s{\cal B}(a,b;s),\\
 {\cal T}_\xi(a,b;0)
 &=3{\cal B}(a,b;0)-2+\xi^2{\cal B}(\xi a,\xi b;0).
\end{align}
Terms multiplied by $\xi$ or $\xi^2$ are taken as their vanishing
limits at $\xi=0$, not as an evaluation of an undefined $0\log0$.

\section{Assembly with the original tadpole prescription}
Let $a_A$ be vector mass squares, $b_i$ the eigenvalues of $H_\xi$,
and $u_i$ its orthonormal eigenvectors. All sums below include ordered
internal pairs. The radial contributions are
\begin{align}
 \Pi^{VV}_{rs}(s)
 &=\frac1{32\pi^2}\left[
  \sum_A C_{rs,AA}{\cal A}^V_\xi(a_A)
  -\sum_{AB}B_{r,AB}B_{s,BA}{\cal T}_\xi(a_A,a_B;s)\right],\\
 \Pi^c_{rs}(s)
 &=\frac1{16\pi^2}\left[
  -\sum_A\xi C_{rs,AA}{\cal A}(\xi a_A)
  +\sum_{AB}\xi^2 B_{r,AB}B_{s,BA}
                         {\cal B}(\xi a_A,\xi a_B;s)\right],\\
 \Pi^{SV}_{rs}(s)
 &=-\frac{4g^2}{16\pi^2}\sum_{A i}
 (d_{Ar}^Tu_i)(d_{As}^Tu_i){\cal I}_\xi(a_A,b_i;s),\\
 \Pi^S_{rs}(s)
 &=\frac1{32\pi^2}\left[
 \sum_i (u_i^TU^\xi_{rs}u_i){\cal A}(b_i)
 -\sum_{ij}(u_i^TT^\xi_ru_j)(u_j^TT^\xi_su_i)
                           {\cal B}(b_i,b_j;s)\right].
\end{align}
The minus sign and factor four in $\Pi^{SV}$ follow from the two
off-diagonal blocks in $\Tr\log{\cal Q}_B$, not from counting a mixed
diagram by analogy. The two $W$ insertions supply $4g^2p_\mu p_\nu$.
In strict Landau gauge the direct scalar-dependent ghost vertices
vanish; this says nothing about ghost contributions to gauge-background
correlators, Yukawa matching or boxes.

With scalar and vector tadpoles $t_S,t_V$, the common finite invariant
mass shift obeys
\begin{equation}
 A_r\,\delta m=-t_S-t_V,\quad
 A_r=-\diag(12w/5,\sigma,v_s),\quad
 H_{\rm CT}=-\diag((t_S+t_V)/r).
\end{equation}
It is the already fixed prescription, not a new zero-momentum
subtraction of each diagram. The Landau radial inverse background kernel is
\begin{equation}
 {\cal K}_B(s)=sG+R^THR+H_{\rm CT}
                      +\Pi^S(s)+\Pi^{VV}(s)+\Pi^{SV}(s).
 \label{eq:total}
\end{equation}
At $s=0$, $\Pi^{SV}=0$ and the vector kernel exactly reproduces the
existing vector CW Hessian with constant $5/6$. The hard scalar CW
comparison omits only the singular soft--soft bubble at $s=0$.
It is \emph{not} an assertion that the full scalar Hessian has a finite
zero-momentum limit. For every $s>0$, all 328 scalar internal states
are included.

The gauge-dependent local wave UV residue is
\begin{equation}
 \Pi_{\rm pole}(s)\big|_{\partial^2}
 =-\frac{g^2(3+\xi)}{16\pi^2\epsilon}
   s\sum_{AI}d_{Ar,I}d_{As,I}.
\end{equation}
The minimal wave counterterm has the opposite sign. No finite wave
shift is introduced. This residue is for the declared background
external-field convention, not a transplanted quantum-field convention.

\subsection{Fermions stay as an explicit input function}
On the already verified radial Yukawa slice $M_F=\sigma F_M$.
If its three Takagi masses are $M_i$, the fixed-VEV Majorana addition is
\begin{equation}
 \Pi^{F,\rm FV}_{\sigma\sigma}(s)
 =\frac1{16\pi^2}\sum_{i=1}^3
    (M_i/\sigma)^2(4M_i^2+s){\cal B}(M_i^2,M_i^2;s),
 \qquad \Pi^{F,\rm FV}_{rs}=0\quad (rs\ne\sigma\sigma).
\end{equation}
This includes its own corresponding fermion tadpole counterterm; do not
also insert that counterterm into $H_{\rm CT}$ in \eqref{eq:total}.
The implementation of this function already exists in
\texttt{verify\_p54\_radial\_fermion\_momentum.py}.
No synthetic masses are promoted to fitted Yukawa data here.

\section{Independent numerical checks and actual result}
\subsection{Checks that do not reuse the master reduction}
For UV-convergent momentum differences we integrate the original
propagators in four dimensions. With $t=k^2$, $z=\cos\theta$,
\begin{equation}
 16\pi^2\int\frac{d^4k}{(2\pi)^4}f(t,z)
 =\int_0^\infty dt\,t\,\langle f\rangle_{S^3},\quad
 q^2=t+s+2\sqrt{ts}\,z.
\end{equation}
The transverse trace is $2+(k\cdot q)^2/(k^2q^2)$, and the mixed
numerator divided by $s$ is
$(1-z^2)/(t+a)+\xi z^2/(t+\xi a)$.
We subtract the $s=0$ vector value or mixed slope before integration.
The finite rational constants cancel in these differences; they are
separately tested with the $d$-dimensional gamma-function tadpole and
bubble, including their poles.

For a massless scalar the angular endpoint $q=0$ converges slowly in
an unsplit angular quadrature. We use the independent exact shell identities
\begin{equation}
 \left\langle\frac1{q^2}\right\rangle=\frac1{\max(t,s)},\qquad
 \left\langle\frac{z^2}{q^2}\right\rangle
       =\frac{t+s}{4\max(t,s)^2}.
\end{equation}
Radial integration is then split at $t=s$. Direct angular integrations
at 320 and 640 nodes approach the same answer. The production master
uses 45-digit arithmetic, not this slower endpoint quadrature.
Another independent UV-finite check is
\begin{equation}
 {\cal T}'_0(a,b;0)
 =-\int_0^\infty dt\left[
 \frac{3bt}{(t+a)(t+b)^3}
       +\frac{3}{4(t+a)(t+b)}\right].
\end{equation}
Integration by parts makes its symmetry in $a,b$ explicit; for $a=b$
it equals $-5/(4a)$. This checks the small-momentum master derivative
without finite-differencing the model Hessian.

\subsection{No new scalar-background or coupling scan}
The radial calculation reuses ten cached polynomial action jets with
zero new full Hessian evaluations. The massive vector squares cluster at
\begin{equation}
 .00506506436,\quad .02532532182,\quad
 .31652201182,\quad .32158707618.
\end{equation}
Different representation sectors can share a mass, so the 35-sector
catalogue is not replaced by the 24 numerical scalar mass clusters.
All multiplicities and off-diagonal vertex contractions are retained.

The complete bosonic radial background kernel \eqref{eq:total} has the
following generalized eigenvalues, solving $\mathcal K_Bz=\lambda Gz$:
\begin{center}
\begin{tabular}{rrrr}
\toprule
$s/\omega^2$&$\lambda_1/\omega^2$&$\lambda_2/\omega^2$&$\lambda_3/\omega^2$\\
\midrule
.001&-.758575&.091032&1.138172\\
.01 &-.726255&.101339&1.730301\\
.05 &-.589137&.142534&2.188925\\
.1  &-.421626&.193296&2.429713\\
.5  & .596222&.795538&3.361713\\
1   &1.098602&2.101012&4.194060\\
\bottomrule
\end{tabular}
\end{center}
These are Euclidean off-shell diagnostics, \emph{not} six particle
mass spectra. Neither a sign change in this table nor a negative
truncated eigenvalue alone establishes a controlled physical pole.
For example, at $s=.05$, adding vector momentum and the mixed bubble
changes the smallest value from $-.584214$ to $-.589137$.

The local gauge insertion, including the mixed term, is
\begin{equation}
 \Delta Z_{\rm gauge}=
 \begin{pmatrix}
 .0008027232&.0027101305&0\\
 .0027101305&-.2951305062&0\\
 0&0&0
 \end{pmatrix}.
\end{equation}
Its generalized eigenvalues are $(-.1475756,0,.0003448)$.
Adding it to the earlier \emph{hard scalar local} insertion gives
\begin{equation}
 \operatorname{eig}_G(\Delta Z_{S,\rm hard}+\Delta Z_{\rm gauge})
       =(.0089952,.2733736,2.4665928).
\end{equation}
This particular diagnostic excludes the nonlocal scalar soft part and
unfitted fermions; it is not the derivative of a finite full kernel at
$s=0$. Nevertheless the omitted gauge contribution does not remove
the earlier large-insertion warning in this convention.

\subsection{Landau limit versus Nielsen certification}
At $s=.05$, keeping the original finite mass CT fixed, the norms
$\|{\cal K}_\xi-{\cal K}_0\|_F$ for
$\xi=.1,.01,.001,.0001$ are respectively
\[
 .9903655,\quad .2646452,\quad .0446276,\quad .00622907.
\]
Smaller $\xi=10^{-5},10^{-6}$ give $.00079812,.00009731$.
Finite-$\xi$ tadpole residuals are recorded, not silently
subtracted by a different parameter choice. Direct ghost terms vanish
in this limit.

There is also an analytic continuity check. At fixed $s>0$, the finite
number of scalar bubble parameter integrals is continuous as a
nonnegative mass approaches zero; the endpoint logarithms are integrable.
Tadpoles $a\log a$ vanish. The explicit ghost terms carry $\xi,\xi^2$,
and the vector/mixed formulas above have the same limiting masters.
Thus the declared functional has a continuous Landau limit, with
mass-log endpoint terms typically of order $\xi|\log\xi|$.
This is not the fixed-bare Nielsen identity: the background displacement,
BRST insertions and background-to-quantum conversion have not been
calculated. Gauge-independent complex poles require that further
structure, not equality of off-shell matrices \cite{GambinoGrassi}.

% Standalone section fragment for inclusion in the same-action momentum note.
% Requires amsmath and amssymb. No independent preamble or bibliography.
\section{Exact tensor-action certification of the transverse projectors}
\label{sec:exact-transverse-action-certificate}

The earlier integer projectors were exact as projectors, but their
identification with the scalar action was only tested numerically.  We now
close that particular gap by an independent Gaussian-integer contraction
calculation.  No floating Hessian, eigensolver, tolerance, or rational fit to
a Hessian enters this certificate.  The claim remains a theorem about the
declared real scalar action on its fixed-orientation radial slice, not a
complete gauged pole or vacuum-stability theorem.

\subsection{A rational background and the canonical five-form map}

Let $q=(q_1<\cdots<q_5)$ label the $252$ five-element subsets of
$\{0,\ldots,9\}$.  Pair $q$ with its complement $q^c$ in the same order
used in the action.  If $s_q$ is the sign of $(q,q^c)$ and the chosen
Hodge eigenvalue is $d\,i$, $d\in\{1,-1\}$, the canonical real
$126+126$ coordinates obey
\begin{align}
 \Sigma_q&=\frac{x_k+i x_{k+126}}2,\\
 \Sigma_{q^c}&=-d\,i\,s_q\frac{x_k+i x_{k+126}}2.
\end{align}
Thus $C_j=2\,\partial\Sigma/\partial x_j$ has Gaussian-integer
quint coefficients.  No $\sqrt2$ remains in this canonical Sigma map.

Choose
\begin{equation}
 (w,\sigma,v_s)=(5,\sqrt2,\sqrt2),\qquad
 \Phi=\operatorname{diag}(-2,-2,-2,-2,-2,-2,3,3,3,3),\quad S=1.
\end{equation}
Writing $B=\Sigma(x_0)$, the tensor $A=4B$ has coefficient
$i^{n_o(q)}$ when $q$ selects one index from each ordered pair
$(0,1),(2,3),\ldots,(8,9)$, and zero otherwise; $n_o(q)$ counts odd
indices.  Hence $A\in\mathbb Z[i]^{252}$.  Antisymmetric extension
uses only permutation signs and zeros for repeated indices.
The certificate also checks the Hodge identities, the real canonical
metric, and the exact background reconstruction.

The output Sigma covectors retain their canonical normalization.
For output Phi covectors we may use $E_{ab}+E_{ba}$, $a<b$, and
$E_{aa}-E_{99}$, $a<9$; for output $H,S$ covectors use $1,i$.
These are invertible real changes of basis, so a vanishing mixed image in
these rational bases is equivalent to a vanishing canonical mixed image.

\subsection{Removing factorial sums before differentiating}

For sorted $I$ define the compressed contractions
\begin{align}
 M_{mn}(U,V)&=\sum_{|I|=4}U_{Im}V_{In},\\
 N_{ab,cd}(U,V)&=\sum_{|I|=3}U_{Iab}V_{Icd},
 \qquad a<b,\ c<d.
\end{align}
An unsorted pair is extended antisymmetrically and a repeated pair is zero.
The original normalization factors $4!$, $3!$, and the ordered-pair factors
then give exactly
\begin{align}
 I_{\lambda_2}&=\langle M(\Sigma,\bar\Sigma),M(\Sigma,\bar\Sigma)\rangle,\\
 I_{\lambda_4}&=\langle N(\Sigma,\bar\Sigma),N(\Sigma,\bar\Sigma)\rangle,\\
 I_{\lambda'_4}&=\langle N(\Sigma,\bar\Sigma),\mathcal C N(\Sigma,\bar\Sigma)\rangle,\\
 (\mathcal C N)_{ab,cd}
 &=N_{ac,bd}-N_{bc,ad}-N_{ad,bc}+N_{bd,ac}.
\end{align}
Here $\langle X,Y\rangle=\sum X_{ij}Y_{ij}$ is a bilinear dot product,
\emph{not} a Hermitian norm.  Confusing the two would erase the important
holomorphic cancellations.  The finite integer permutation matrix
$\mathcal C$ is checked to be self-adjoint for this bilinear dot product.

For either $X=M$ or $X=N$, define
\begin{align}
 \widehat X_0&=X(A,\bar A),\\
 G_j&=X(C_j,\bar A)+X(A,\bar C_j),\\
 Q_{jk}&=X(C_j,\bar C_k)+X(C_k,\bar C_j).
\end{align}
Since $B=A/4$ and $\partial_j\Sigma=C_j/2$, differentiating the
quadratic Gram functional gives the exact Gaussian-integer formula
\begin{equation}
 \boxed{32(H_\lambda)_{jk}
 =\langle G_j,\mathcal C G_k\rangle
  +\langle\widehat X_0,\mathcal C Q_{jk}\rangle.}
\end{equation}
For $\lambda_2,\lambda_4$, replace $\mathcal C$ by the identity.
The resulting imaginary parts vanish exactly.  The matrices are at most
$252\times252$; the compressed Gram arrays are $10\times10$ and
$45\times45$.

\subsection{The mixed covectors are part of the certificate}

The other nontrivial contractions with a Sigma input are
\begin{align}
 I_\beta&=\sum_{ijkl}\Phi_{ij}\Phi_{kl}
                  N_{ik,jl}(\Sigma,\bar\Sigma),\\
 I_{\chi_4}&=S\sum_{mn}\Phi_{mn}M_{mn}(\Sigma,\Sigma)+\mathrm{c.c.},\\
 I_{\eta_1}&=\frac16\sum_{a<b,c<d,n}
         N_{ab,cd}(\Sigma,\bar\Sigma)\Sigma_{abcdn}H_n
         +\mathrm{c.c.}
\end{align}
They are differentiated before setting the external $H$ background to zero.
This includes all Phi--Sigma covectors for $\beta$, all Phi--Sigma and
S--Sigma covectors for $\chi_4$, and all twenty H--Sigma covectors for
$\eta_1$.  In particular, the fact that $I_{\eta_1}=0$ at $H=0$ does
not justify dropping its mixed Hessian.
The remaining active norm terms are $\nu^2,\lambda_0,\alpha,\chi_2$.
All nineteen other invariants either have no Sigma input or retain a
factor of $H$ after the two derivatives, and therefore give zero.

\subsection{Finite exact certificates and their continuation}

The stored integer matrices $K_{50},K_{36}$, each $252\times252$,
are treated as candidate matrices, with $P_r=K_r/8$ embedded into the
Sigma block.  The following are checked using exact integer arithmetic:
\begin{equation}
 K_r^T=K_r,\qquad K_r^2=8K_r,\qquad
 \operatorname{tr}K_r=8r,\qquad K_{50}K_{36}=0.
\end{equation}
Symmetry and idempotency make the trace an exact rank certificate.
For every one of the $29$ unit invariants, its independently derived
rational $328\times252$ covector matrix $H_\alpha=N_\alpha/d_\alpha$
is checked against
\begin{equation}
 N_\alpha K_r=d_\alpha c_{\alpha,r}\,\iota_\Sigma K_r,
 \qquad r=50,36.
\end{equation}
All denominators are cleared; every residual entry is the integer zero.
At the rational background the only nonzero eigenvalues are
\begin{equation}
\begin{array}{c|rrrrrrrr}
 \text{invariant}&\nu^2&\lambda_0&\lambda_2&\lambda_4&\lambda'_4&\alpha&\beta&\chi_2\\ \hline
 c_{\alpha,50}&-\frac12&1&8&8&32&30&-30&120\\
 c_{\alpha,36}&-\frac12&1&12&12&16&30&-30&120.
\end{array}
\end{equation}
Every individual invariant Hessian block has a single radial multidegree:
the invariant degree minus the degrees of its two external legs.
Thus each zero block and scalar Sigma block proved at this nonzero radial
point extends exactly to every nonzero radial point of the same angular
orientation.  Linearity in the real action parameters then proves
\begin{align}
 m_{50}^2={}&-\frac{\nu^2}{2}
 +\left(\frac{\lambda_0}{2}+4\lambda_2+4\lambda_4+16\lambda'_4\right)\sigma^2
 +\frac65(\alpha-\beta)w^2+60\chi_2v_s^2,\\
 m_{36}^2={}&-\frac{\nu^2}{2}
 +\left(\frac{\lambda_0}{2}+6\lambda_2+6\lambda_4+8\lambda'_4\right)\sigma^2
 +\frac65(\alpha-\beta)w^2+60\chi_2v_s^2.
\end{align}
Consequently the action-to-projector assumption in the earlier
$1/\tau$ argument has been removed.  The condition that the retained hard
spectrum is positive, and the restriction to the scalar hard one-loop
kinetic kernel, are not removed.

\paragraph{Arithmetic audit.}
The implementation uses pairs of signed integer arrays for Gaussian
integers. Before each large contraction product and sum it bounds every
resulting entry by the corresponding dimension times input entry bounds.
The fixed geometry constructors separately use only $0,\pm1$, the
diagonal entries $-2,3$, and explicitly small integer factors.
The largest recorded conservative bound is $1{,}008{,}000$, safely below
$2^{62}$ and hence below signed integer overflow.  No small residual is
accepted: the certificate consists of exact zero-entry tests.  The source
action is pinned by its SHA-256 digest; a changed action requires a fresh
audit of the symbolic contraction transcription.


\section{What the exact identities do and do not repair}
With only $(\lambda_2,\lambda_4,\lambda'_4)$ multiplied by $\tau>0$
and the stationary mass relation imposed, the now-certified blocks give
\begin{equation}
 m_{50}^2=12\tau\sigma^2,\qquad
 m_{36}^2=\frac{58}{5}\tau\sigma^2.
\end{equation}
Fluctuation derivatives must be taken \emph{before} imposing that
stationary relation. For $c_\kappa=12,58/5$ they are
\begin{equation}
 b_\kappa=\left(\tfrac{12}{5}(\alpha-\beta)w,\,
       (\lambda_0+2c_\kappa\tau)\sigma,\,
       120\chi_2v_s\right),\quad b_0=(.36,.11385,.03).
\end{equation}
For a positive retained hard spectrum, the scalar pair kernel has
positive weights. Keeping only these two diagonal blocks gives
\begin{equation}
 \Delta Z_{S,\rm hard}\succeq
 \sum_{\kappa}\frac{N_\kappa b_\kappa b_\kappa^T}
                       {192\pi^2c_\kappa\tau\sigma^2}.
\end{equation}
Use the fixed $G$-unit witness
$z_0=G^{-1}b_0/\sqrt{b_0^TG^{-1}b_0}$.
Since $b_\kappa^Tz_0\ge b_0^Tz_0>0$, one obtains the exact
coefficient formula
\begin{equation}
 \lambda_{\max}(G^{-1/2}\Delta Z_{S,\rm hard}G^{-1/2})
 \ge\frac{b_0^TG^{-1}b_0}{192\pi^2\sigma^2\,\tau}
      \left(\frac{50}{12}+\frac{36}{58/5}\right)
 =\frac{.01471609503\ldots}{\tau}.
\end{equation}
The decimal is a presentation of the preceding exact expression,
not a fitted constant. No false matrix inequality
$b_\kappa b_\kappa^T\succeq b_0b_0^T$ is used. The exact certificate
removes the action-to-projector assumption in this argument. It does
not turn a positive scalar subset into a bound on a signed full gauge
or fermion kernel, nor into a whole-model no-go theorem.

\section{Next decisions and reproducibility}
The two requested subgates have changed status: the explicit radial
bosonic momentum functional is evaluated, and the real-action transverse
identities have an exact certificate. Default quartics, the old vacuum,
matching trajectories and observed-mass targets were not altered.

The next bounded options are:
\begin{enumerate}
 \item Derive the fixed-bare Nielsen/BRST insertions and the complete
 neutral physical mixing block, including the background/quantum map.
 Track the induced background displacement and light-doublet retuning
 in the same tadpole prescription before interpreting a pole.
 \item Use the certified vertex-to-gap criterion for an explicit
 retained-light-block EFT and momentum power counting. Do not integrate
 these blocks out merely because their tree masses are nonzero.
 Any new candidate needs mixed-field boundedness and competing-vacuum
 tests as well as perturbative remainder control.
\end{enumerate}
Neither route authorizes another blind quartic scan. Full Wilson/box,
lower gauge/ghost/Yukawa finite matching and the complete finite seesaw
including $C_{HN}$ remain open. The one-loop matrix alone cannot solve
these other gates. Goldstone resummation methods address specific IR
problems, not automatically this complete gauged matching problem
\cite{BraathenGoodsell}.

The new verifiers are
\begin{quote}\small
\path{route_f/code/verify_p54_gauge_mixed_momentum.py}\\
\path{route_f/code/verify_p54_exact_transverse_certificate.py}.
\end{quote}
The gauge verifier passes 99/99 numerical/regression checks; the action
certificate passes 109/109 exact checks. Their JSON files retain
matrices, individual checks and source hashes, with Markdown summaries.
The gauge verifier uses the existing P1 JAX requirements plus NumPy,
SciPy and mpmath; the exact certificate uses only Python, NumPy and SciPy,
not JAX or any floating Hessian cache. Run the former with the existing
\texttt{--cache-dir}; the latter needs no arguments. Exact proof claims
refer to zero integer residuals for the pinned action, whereas quadrature
and floating matrix checks retain their explicit numerical tolerances.

\begin{thebibliography}{9}
\bibitem{BraathenGoodsell}
J.~Braathen and M.~D.~Goodsell,
\emph{Avoiding the Goldstone Boson Catastrophe in general renormalisable
field theories at two loops},
\href{https://arxiv.org/abs/1609.06977}{arXiv:1609.06977}.
\bibitem{GambinoGrassi}
P.~Gambino and P.~A.~Grassi,
\emph{The Nielsen Identities of the SM and the definition of mass},
\href{https://arxiv.org/abs/hep-ph/9907254}{arXiv:hep-ph/9907254}.
\end{thebibliography}
\end{document}
